\documentclass{thesisclass}

%% -------------------------------
%% |    IM Thesis Template       |
%% -------------------------------
%% Further additions by: Julian Huber, 2019
%% julian.huber "at" kit.edu
%% Further additions by: Kenneth Qua, 2021
%% kenneth.qua "at" wu.ac.at

%% Notes:
%% Language switch after \begin{document}

% Based on thesisclass.cls of Timo Rohrberg, 2009
% ----------------------------------------------------------------
% Thesis - Main document
% ----------------------------------------------------------------

% Select your thesis language here:
\usepackage[english]{babel}
%\usepackage[german]{babel}


%% ---------------------------------
%% |      Additional packages      |
%% ---------------------------------
%% 
\usepackage{algorithm}
\usepackage{listings}
\usepackage{algpseudocode}
\usepackage{amssymb} 

\usepackage{graphicx}
%http://en.wikibooks.org/wiki/LaTeX/Importing_Graphics#Graphics_storage
\DeclareGraphicsExtensions{.pdf,.png,.jpg}
\graphicspath{{./figures/}} %Use curly braces for each path to add and don't
% forget trailing slash '/'
\providecommand{\ResGraphic}[2]{\IfFileExists{figures/#2.pdf}{\includegraphics[width=#1\linewidth]{#2}}{\framebox[#1\linewidth][c]{\rule[-1.1cm]{0pt}{2.4cm}\emph{figure pending}}}}
\providecommand{\ResPendingTable}[2]{\begin{table}[ht]\centering\caption{#1}\label{#2}\emph{To be completed from the full run.}\end{table}}
% \usepackage{epstopdf} %Nice to automatically convert eps figures to pdf
% format  (from inkscape, etc)
\usepackage[natbibapa]{apacite}
\usepackage{booktabs}

%% ---------------------------------
%% | Needed for the List of Abbreviations |
%% ---------------------------------
\usepackage{nomencl}
\renewcommand{\nomname}{List of Abbreviations}
% Punkte zw. Abkürzung und Erklärung
\setlength{\nomlabelwidth}{.32\hsize}%.32\hsize bezieht sich auf die Punkte zwischen Abkürzung und Erklärung -> je größer die Zahl, desto mehr %Punkte
\renewcommand{\nomlabel}[1]{#1 \dotfill}
% Zeilenabstände verkleinern
\setlength{\nomitemsep}{-\parsep}
\makenomenclature
\usepackage[normalem]{ulem}
\newcommand{\markup}[1]{\uline{#1}}

\usepackage{xspace}

\newcommand\ie{i.\,e.\xspace}
\newcommand\eg{e.\,g.\xspace}
\newcommand\Eg{E.\,g.\xspace}
\newcommand{\Latex}{\LaTeX\xspace}


\usepackage{tikz}
\usetikzlibrary{arrows.meta, positioning, bending}
\definecolor{scamPurple}{HTML}{7F77DD}
\definecolor{scamTeal}{HTML}{1D9E75}
\definecolor{scamCoral}{HTML}{D85A30}
\definecolor{scamBlue}{HTML}{378ADD}
\tikzset{
scamnode/.style={circle, minimum size=11mm, inner sep=0pt},
scamedge/.style={-{Stealth[length=2.4mm,width=2.2mm]}, line width=0.9pt},
scamedgeheavy/.style={-{Stealth[length=2.8mm,width=2.8mm]}, line width=1.8pt},
edgelbl/.style={inner sep=2pt, font=\small},
every label/.append style={font=\small}
}

%% ---------------------------------
%% |   Code listing style (lstlistings)   |
%% ---------------------------------
\definecolor{lstbg}{HTML}{F7F7F9}
\definecolor{lstframe}{HTML}{D0D0D8}
\definecolor{lstkw}{HTML}{1D4E89}
\definecolor{lstcomment}{HTML}{4B8B3B}
\definecolor{lststring}{HTML}{A8431E}
\definecolor{lstnum}{HTML}{9A9A9A}
\lstset{
  basicstyle=\ttfamily\footnotesize,
  backgroundcolor=\color{lstbg},
  keywordstyle=\color{lstkw}\bfseries,
  commentstyle=\color{lstcomment}\itshape,
  stringstyle=\color{lststring},
  numberstyle=\tiny\color{lstnum},
  numbers=left,
  numbersep=8pt,
  frame=single,
  rulecolor=\color{lstframe},
  breaklines=true,
  breakatwhitespace=true,
  postbreak=\mbox{\textcolor{lstnum}{$\hookrightarrow$}\space},
  columns=fullflexible,
  keepspaces=true,
  showstringspaces=false,
  tabsize=2,
  captionpos=b,
  xleftmargin=14pt,
  framexleftmargin=14pt,
  aboveskip=1.2em,
  belowskip=1em,
}



%% ---------------------------------
%% | Information about the thesis  |
%% ---------------------------------

\newcommand{\mytype}{\iflanguage{english}{Bachelor Thesis}{Bachelorarbeit}} 

\newcommand{\myname}{Sebastian Eduard Bruckner-Hrubesch}
\newcommand{\matricle}{h12220272}
\newcommand{\mytitle}{\iflanguage{english}{Coordinated Malicious Behavior Detection in Blockchain Networks Using Temporal Graph Motifs}{Titel der Arbeit}}
\newcommand{\myinstitute}{\iflanguage{english}
{Institute for Complex Networks}
{Institut für Komplexe Netzwerke}}

\newcommand{\advisor}{Mag.\ Ema Kahr, PhD}

\newcommand{\timestart}{\iflanguage{english}{03. April 2026}{03. April 2026}}
\newcommand{\timeend}{\iflanguage{english}{XXth of Month 20XX}{XX. Monat 20XX}}
\newcommand{\submissiontime}{DD. MM. 20XX}

%% -------------------------------
%% |  Information for PDF file   |
%% -------------------------------
%% IM: Auto-Fill this information
\hypersetup{
 pdfauthor={\myname},
 pdftitle={\mytitle},
 pdfsubject={\mytype},
 pdfkeywords={\mytype}
}

%% ---------------------------------
%% | ToDo Marker - only for draft! |
%% ---------------------------------
% Remove this section for final version!
%\setlength{\marginparwidth}{20mm}

%\newcommand{\margtodo}
%{\marginpar{\textbf{\textcolor{red}{ToDo}}}{}}

%\newcommand{\todo}[1]
%{{\textbf{\textcolor{red}{(\margtodo{}#1)}}}{}}


%% --------------------------------
%% | Old Marker - only for draft! |
%% --------------------------------
% Remove this section for final version!
%\newenvironment{deprecated}
%{\begin{color}{gray}}
%{\end{color}}


%% --------------------------------
%% | Settings for word separation |
%% --------------------------------
% Help for separation:
% In german package the following hints are additionally available:
% "- = Additional separation
% "| = Suppress ligation and possible separation (e.g. Schaf"|fell)
% "~ = Hyphenation without separation (e.g. bergauf und "~ab)
% "= = Hyphenation with separation before and after
% "" = Separation without a hyphenation (e.g. und/""oder)

% Describe separation hints here:
\hyphenation{
% Pro-to-koll-in-stan-zen
% Ma-na-ge-ment  Netz-werk-ele-men-ten
% Netz-werk Netz-werk-re-ser-vie-rung
% Netz-werk-adap-ter Fein-ju-stier-ung
% Da-ten-strom-spe-zi-fi-ka-tion Pa-ket-rumpf
% Kon-troll-in-stanz
}


%% ------------------------
%% |    Including files   |
%% ------------------------
% Only files listed here will be included!
% Userful command for partially translating the document (for bug-fixing e.g.)
\includeonly{%
titlepage,
text/abbreviations,
text/introduction,
text/content,
text/evaluation,
text/conclusion,
text/appendix
}

%%%%%%%%%%%%%%%%%%%%%%%%%%%%%%%%%
%% Here, main documents begins %%
%%%%%%%%%%%%%%%%%%%%%%%%%%%%%%%%%
\begin{document}

\frontmatter
\pagenumbering{roman}
\include{titlepage}
% IM Style: No additional blank page
% \blankpage


%% -------------------
%% |   Directories   |
%% -------------------
\iflanguage{english}{
\renewcommand*\contentsname{Table of Contents}}
{}
\clearpage
\tableofcontents

% IM Style: No additional blank page
% \blankpaage

% If you do not wish to include a list of figures, list of tables and list of abbreviations, comment out the following section
\iflanguage{english}
{
\clearpage
\phantomsection % Properly links List of Figures in Table of Contents and starts it on new page
\listoffigures \addcontentsline{toc}{chapter}{List of Figures} 
\clearpage
\phantomsection % Properly links List of Tables in Table of Contents and starts it on new page
\listoftables  \addcontentsline{toc}{chapter}{List of Tables} 
\printnomenclature   \addcontentsline{toc}{chapter}{List of Abbreviations} 
}
{
\clearpage
\phantomsection % Properly links List of Figures in Table of Contents and starts it on new page
\listoffigures \addcontentsline{toc}{chapter}{Abbildungsverzeichnis} 
\clearpage
\phantomsection % Properly links List of Tables in Table of Contents and starts it on new page
\listoftables \addcontentsline{toc}{chapter}{Tabellenverzeichnis} 
\printnomenclature \addcontentsline{toc}{chapter}{Abkürzungsverzeichnis}
}


% Execute this command for index creation, i.e., for abbreviations by the nomencl package
% makeindex thesis.nlo -s nomencl.ist -o thesis.nls



%\printnomenclature



%% -----------------
%% |   Main part   |
%% -----------------
\mainmatter
\pagenumbering{arabic}


%% ==============
\chapter*{Abstract} 
This section is a short summary of your entire paper. It should briefly cover the most important aspects of your paper. After reading this section, the reader should have a good understanding of what the research question is, how you answered it and what the answer was. This section is intentionally meant to be brief and no longer than half a page.


Notes:
* how to formalize coordinated behavior.
* anti-money laundering typology -> transfer to dataset -> network motive (typical sequence)
* how to compare if it is statistically significant 
* encode search of temporal motive 

Look into related work
Dig into temporal motives
Isomorphism testing -> no, but counting
Not Motive mining but “counting”
* timestamps

networkX -> python toolkit
Create network from this toolkit

Pool of known motives


\chapter{Introduction}
Blockchain networks have become a critical part of the global financial infrastructure. As they have grown, they have also become an attractive target for criminal actors. According to the Chainalysis 2025 Crypto Crime Report, more than \$40.9 billion flowed to illicit cryptocurrency addresses in 2024 alone \citep{chainalysis2025, shi2025realcats}. This number is likely a lower bound. Chainalysis notes that annual illicit-volume estimates have grown by an average of 25\% between reporting periods since 2020, so the 2024 total could exceed \$51 billion \citep{chainalysis2025}. The abuse ranges from ransomware and phishing scams to large-scale money laundering, and it affects both the Bitcoin and the Ethereum ecosystem \citep{shi2025realcats, sanjalawe2023atd}.

The harder problem is not the size of these losses, but the way they are produced. Sophisticated attackers do not rely on a single, continuously active address. Instead, they spread their activity over many short-lived wallets, where each wallet makes only a few transactions that look unremarkable on their own. In fact, 90.70\% of criminal Bitcoin addresses carry a zero balance and 60\% of them stop being active within 31 days \citep{shi2025realcats}. This behaviour deliberately imitates ordinary wallets, so a single transaction rarely looks suspicious. The signal of coordinated activity therefore does not appear in individual transactions, but in how those transactions are connected across several addresses and over time. A ransomware payment chain, a phishing fan-out or a layering operation only becomes visible once its transactions are considered together, in order, and across all the addresses involved.

Detecting these patterns requires a method that explicitly models the temporal order of transactions. However, most existing approaches do not do this. Feature-engineering methods extract aggregate per-address statistics such as transaction count, address lifetime and balance variance \citep{chen2022phishing, alemari2021labeled}. Such features are easy to evade by spreading activity over many short-lived wallets. Graph-embedding methods capture the structure of the transaction graph, but not the temporal order of events across several addresses \citep{shi2025realcats}. Both approaches share the same limitation: they cannot capture ordered, multi-party interaction patterns that span several addresses and several time steps, and this is exactly where coordinated malicious behaviour shows itself.

Temporal motifs are a natural way to close this gap. A temporal motif is a recurring, time-ordered subgraph pattern, where a set of edges occurs among a group of nodes within a bounded time window \citep{kovanen2011temporal, paranjape2017motifs}. Applied to the Bitcoin transaction graph, temporal motif counts already reveal structural anomalies that static analysis cannot detect. For example, cyclic triangle motifs are far more frequent in Bitcoin than in any communication network, which reflects the recirculation that is typical of mixing and layering \citep{paranjape2017motifs}. Despite this, using temporal motifs as features for the supervised detection of coordinated malicious behaviour across multiple addresses is still largely unexplored. This thesis addresses that gap directly.

TODO: research question(s)

This thesis investigates the following research questions:

\begin{description}
\item[RQ1.] Do temporal motifs improve the detection of malicious Ethereum addresses compared to transaction aggregates and static graph features alone?
\item[RQ2.] Which temporal window size (\delta) maximizes the usefulness of temporal motif fingerprints for detecting coordinated malicious behaviour?
\item[RQ3.] Do different categories of coordinated malicious behaviour exhibit distinct temporal motif fingerprints?
\item[RQ4.] Does the temporal window that best captures these fingerprints differ across behaviour categories?
\end{description}

\chapter{Related Work} 

how the other papers analyze blockchain in terms of malicious behaviour

dont go into scamming in general -> only network focused

this and that paper mentioned....

- 2-3 pages

\section{Graph-Based Detection of Malicious Behavior on Blockchain Networks}

The most common approach to detecting malicious behavior on blockchain networks is
to model the transaction history as a directed graph and extract features from it.
Sanjalawe and Al-Emari~\cite{sanjalawe2023atd} represent the Ethereum network
as a weighted multi-digraph and collect features such as degree statistics,
average transaction amounts, and gas metadata, with a 95.06\% detection accuracy.
Chen et al.~\cite{chen2022phishing} take a similar approach for phishing detection,
constructing 149 address-level features inspired by fraud patterns in traditional
bank accounts and reaching 98.89\% precision with a Random Forest classifier.

However those approaches either do not take the temporal ordering of
transactions into consideration or reduce it to aggregated statistics like the address lifetime.
Coordinated malicious behavior shows itself in the order in which transactions
occur across multiple addresses within a specific time window, a pattern that
static node features and embeddings cannot capture.
~\cite{paranjape2017motifs} demonstrate this directly on Bitcoin:
the fraction of cyclic triangle motifs is substantially higher than in any
communication network they study, a property that only becomes visible when the
temporal ordering of edges is kept.
This makes a detection approach based on temporal graph motifs rather than
static graph representations necessary.

\section{Datasets for Blockchain Transaction Analysis}
\label{sec:relwork_datasets}

Two publicly available datasets are central to the related literature reviewed in this
thesis: the Benchmark Labelled Transactions Ethereum (BLTE) dataset, which is used
as the primary evaluation benchmark in the intrusion detection work of
Sanjalawe and Al-Emari~\cite{sanjalawe2023atd}, and the Real-CATS dataset,
which is introduced as a cross-chain training resource for general cryptocurrency
cybercrime detection~\cite{shi2025realcats}.

\subsection*{BLTE}

The BLTE dataset was originally proposed by Al-Emari et al.~\cite{alemari2021labeled}
as a labeled transactions-based benchmark for the Ethereum network.
It is built on the Ethereum Classic (ETC) chain and is publicly available on
Kaggle.
Fraud addresses are sourced from EtherScamDB, an open-source registry of scam
addresses, from which 2{,}357 unique Ethereum addresses were extracted,
including phishing and scamming categories.
The corresponding transactions were retrieved for each address over the period from
2017-2019, with a total of 262{,}123{,}807 raw transactions of which
16{,}784 are marked abnormal~\cite{alemari2021labeled}.
The final dataset contains 71{,}250 transactions
at a 4:1 normal-to-abnormal ratio, with each transaction represented by 18 attributes.
~\cite{alemari2021labeled}.

The BLTE dataset serves as the direct evaluation environment for the
approach of Sanjalawe and Al-Emari~\cite{sanjalawe2023atd}.
Their method models the Ethereum network as a weighted multi-digraph
\(G = (N, E)\) with edge tuples \((s, u, v, t)\) encoding sender, receiver,
value, and timestamp, from which 22 node-level features are derived via
multi-digraph theory~\cite{sanjalawe2023atd}.
A feature selection procedure then reduces this set to
7 mutual features, including degree statistics, average transaction amounts,
active duration, mean inter-transaction interval, and average gas
price~\cite{sanjalawe2023atd}.
A Semi-Supervised Generative Adversarial Network classifier trained on these
features achieves a detection accuracy of 95.06\% and an F1-measure of 95.11\%
on the BLTE data, outperforming standard baselines such as Random Forest, SVM,
and LSTM~\cite{sanjalawe2023atd}.

\subsection*{Real-CATS}

The Real-CATS dataset was proposed by Shi et al.~\cite{shi2025realcats} to address
the generalizability and comparability shortcomings of existing blockchain datasets.
Real-CATS is the largest publicly available labeled dataset for cryptocurrency
cybercrime detection, comprising 103{,}203 criminal addresses and 106{,}196 benign
addresses across Bitcoin and Ethereum.
Criminal addresses are derived from real-world abuse reports crawled from
Chainabuse.com, while benign addresses are collected from the outputs of exchange
hot wallets based on the assumption that centralized exchanges apply KYC
procedures~\cite{shi2025realcats}.
Concretely, the Bitcoin split contains 90{,}612 criminal and 90{,}176 benign
addresses (CB and BB), and the Ethereum split contains 12{,}591 criminal and
16{,}020 benign addresses (CE and BE), each with a profile of 32 and 52 statistical
features respectively~\cite{shi2025realcats}.
A supplementary testing dataset, Sup-CATS, was additionally collected from all
blocks on a single day (12 August 2024) on Ethereum, comprising 3{,}147
Etherscan-labeled criminal addresses and 11{,}058 benign addresses, to simulate a
real-world deployment scenario~\cite{shi2025realcats}.

The dataset is explicitly designed to satisfy four characteristics: Comprehensiveness,
Classifiability, Customizability, and Real-world Transferability (C3R).
The authors argue that these characteristics are fundamental for practical cybercrime detection
research~\cite{shi2025realcats}.
In terms of customizability, Shi et al.~\cite{shi2025realcats} demonstrate that
the detailed transaction records released alongside the profiles allow the
construction of weighted multi-digraphs, on which the trans2vec random walk
method~\cite{wu2020phishers} achieves 0.90 precision, 0.93 recall, and 0.91
F1-score on a 2{,}500-address subset.
However, Shi et al.~\cite{shi2025realcats} also note that node embedding methods
are computationally expensive and require building the full transaction graph at training
time, making real-time detection of new addresses challenging.

The two datasets have important roles in the related literature.
BLTE provides a transaction-level benchmark for the Ethereum network that was
the first publicly available labeled dataset of this type, but its labeling is
limited to the EtherScamDB registry and its transaction volume is constrained
to the 2017-2019 period~\cite{alemari2021labeled}.
Real-CATS addresses the cross-chain scope and real-world requirements that
BLTE does not cover, offering substantially more addresses and a more diverse
behavioral distribution~\cite{shi2025realcats}.

Both datasets, however, use individual-address aggregates or embeddings in
their published detection baselines. Real-CATS does release detailed
per-address transaction histories, but neither benchmark study evaluates
per-address $\delta$-temporal motif fingerprints with held-out fraud labels.
That untested representation is the focus of this thesis.

\section*{Temporal Motifs in Financial and Blockchain Networks}

A small body of work has begun to apply temporal motif counting to
financial transaction networks, though none of it combines all three
elements that this thesis requires: blockchain data, labeled malicious
behavior, and temporal motif-based detection.

Arnold et al.~\cite{arnold2024motifs} apply the temporal motif framework
of Paranjape et al.\ directly to two Bitcoin transaction datasets and one
NFT dataset, analyzing sequences of three transactions across up to three
users. They show that the motif distribution is heavy-tailed across users
and that examining per-user motif signatures and motif counts over time
reveals anomalous activity that a global count over all users obscures.
The authors explicitly suggest that local temporal motif profiles could
serve as features for classification or clustering in anti-money
laundering or fraud detection~\cite{arnold2024motifs}. However, the
paper stops short of constructing any detection system: no labeled
criminal addresses are used, and the anomalous activity is identified
only descriptively.

Liu et al.~\cite{liu2023temporal} demonstrate that temporal motifs are
effective features for fraud detection on financial transaction networks,
outperforming static graph features and node embedding baselines such as
LINE and node2vec on labeled datasets from Mercari and a synthetic
J.P.\ Morgan Chase network. The datasets are not blockchain-based,
however, and the work does not address the specific behavioral patterns
of coordinated on-chain scam activity.

Mo et al.~\cite{mo2022motifgcn} incorporate motif matrices into a
temporal graph convolutional network for fraud detection on the
Bitcoin-Alpha and Bitcoin-OTC trust networks. The model captures local
topological structure through motif matrices and temporal evolution
through snapshot-based aggregation. This is the closest existing work to
the framing of this thesis, yet two important distinctions apply: the
motifs used are static subgraph counts computed per snapshot rather than
$\delta$-temporal motifs in the sense of Paranjape et al., and the
datasets are signed trust networks rather than transaction graphs with
labeled scam addresses.

These three papers show that temporal motifs carry signal for financial fraud and that blockchain
transaction data is well suited to temporal motif analysis. However, the
literature does not yet contain a method that counts $\delta$-temporal
motifs on a labeled blockchain scam dataset and uses those motifs directly as features for detection. This combination is the
core contribution of this thesis.

Sources: 
- Google Scholar
- Science Direct
- ACM digital library
- IEEE Xplore digital library
- Wiley
- Springer Link





\chapter{Theoretical Foundations}
This chapter will cover the theoretical foundations that this thesis builds upon.

\section{Blockchain and Ethereum fundamentals}
First, we will cover what a Blockchain is and how Ethereum works and handles transactions. 
A blockchain is a distributed ledger in which transactions are grouped into blocks and linked chronologically so that later blocks reference earlier ones through cryptographic hashes \cite{nakamoto2008bitcoin, buterin2014ethereum}. This makes the ledger append-only, because changing a past record would require rewriting the block and all of the blocks that follow under the network’s consensus rules \cite{nakamoto2008bitcoin}. For this thesis, the important part is not the consensus mechanism itself, but that blockchain data is ordered, time-stamped, and persistent.

Ethereum extends this basic blockchain idea with account-based state and smart-contract execution \cite{buterin2014ethereum}. The relevant transaction record can be represented as a tuple $(s,r,v,t,g)$, where $s$ is the sender, $r$ the receiver, $v$ the transferred value, $t$ the timestamp, and $g$ the gas usage (transaction fee) \cite{alemari2021labeled, sanjalawe2023atd}. Ethereum distinguishes between externally owned accounts (EOAs), which are controlled by private keys and initiate transactions directly, and contract accounts, which execute code when invoked \cite{buterin2014ethereum, alemari2021labeled}.

Because multiple transactions can occur between the same pair of addresses at different times, with different values and gas costs, Ethereum is naturally modeled as a weighted temporal directed multigraph \cite{kovanen2011temporal, paranjape2017motifs}. This is the abstraction used in the remainder of the chapter. In contrast, Bitcoin follows a UTXO (unspent transaction output) model, where transactions consume and create unspent outputs rather than directly transferring value between persistent accounts \cite{nakamoto2008bitcoin}.

With this data model established, the next section introduces the formal graph representation.

\section{Graph-theory foundations}
\label{sec:graph}
what is a graph

A graph is a mathematical structure consisting of a set of nodes $V$ and a set 
of edges $E$, where nodes represent entities and edges represent pairwise 
relations between them \cite{newman2010networks}. In the context of this 
thesis, nodes correspond to Ethereum addresses and edges correspond to 
transactions between them.

Because transactions have a clearly defined sender and receiver, the 
appropriate structure is a \emph{directed graph}, where an edge $(u, v) \in E$ 
represents a directed relation from node $u$ to node $v$ 
\cite{newman2010networks}. Since the same pair of addresses may transact 
multiple times, each with a distinct value and timestamp, the data is more 
precisely represented as a \emph{directed multigraph}, in which multiple edges 
between the same ordered pair of nodes are permitted 
\cite{paranjape2017motifs, sanjalawe2023atd}.

For a node $v$, the \emph{in-degree} $d^-(v)$ counts the number of incoming 
edges and the \emph{out-degree} $d^+(v)$ counts the number of outgoing edges 
\cite{newman2010networks}. These degree measures serve as basic node-level 
features and are used as baselines in the classification experiments in 
Chapter~\ref{ch:method}.

The directed multigraph provides the structural foundation for the graph model 
used throughout this thesis. The following section extends this static 
representation by incorporating the timestamps of transactions, which is 
necessary to capture the temporal ordering of interactions.


\section{Temporal networks}
\label{sec:temporal_networks}
The graph model introduced in Section~\ref{sec:graph} treats the Ethereum transaction graph as a static object once timestamps are discarded. This section explains why that is insufficient for detecting coordinated malicious behavior.

\paragraph{From static to temporal.}
Static networks aggregate all interactions into a single, time-independent structure. While such aggregation is computationally efficient, it completely ignores the ordering of events, which is exactly the information that distinguishes coordinated sequential behavior from coincidental events. Temporal networks take this information into account by retaining the full sequence of timestamped edges \cite{kovanen2011temporal, paranjape2017motifs}.

\paragraph{Time-respecting paths and edge ordering.}
blabla


\paragraph{The time window $\delta$.}
*sliding window
*why delta was chosen

Not all time-respecting paths are meaningful for detecting coordinated behaviour. For example, two transactions that happened with years in between are unlikely to be part of the same coordinated action. To focus on interactions that are \emph{temporally close}, a time window $\delta$ is introduced. Two edges $e_i$ and $e_j$ are considered temporally close if they share at least one endpoint and $|t_j - t_i| \leq \delta$ \cite{kovanen2011temporal, paranjape2017motifs}. Choosing the appropriate $\delta$ is important for our analysis. A small $\delta$ captures rapid bursts of activity such as phishing campaigns that distribute funds within minutes, while a larger $\delta$ captures slower coordination patterns such as layering across multiple days. In this thesis, $\delta$ is treated as a configurable parameter whose effect on motif counts is evaluated in Chapter~\ref{ch:results}. % was: Chapter TODO:blabla

\paragraph{What static projection loses.}
The aggregated static graph loses two things: the \emph{ordering} of edges and the \emph{temporal density} of interactions. This means that two structurally identical subgraphs in the static graph could correspond to entirely different behaviours depending on whether the underlying transactions occurred rapidly or were spread across a bigger time span. \cite{kovanen2011temporal} observe that static networks are often time-aggregates of systems where connections are not continuously active, and that temporal inhomogeneities, which static representations do not capture, have an important effect on dynamic processes. In the Ethereum context, a fan-out pattern in the static graph (one address sending to many others) looks the same regardless of whether all transfers happened within ten seconds or over ten months. Only the temporal graph reveals whether the fan-out was a coordinated event.

\paragraph{Motivation for temporal motifs.}
The observation that edge ordering matters and that meaningful patterns must occur within a bounded time window motivates the concept of temporal motifs introduced in Section~\ref{sec:motifs}. Rather than asking which subgraph topologies are over-represented in the static graph, temporal motifs ask which \emph{ordered sequences} of interactions recur within a window $\delta$. This is the appropriate question for detecting coordinated behavior in blockchain networks, where the timing of transactions is as informative as their structure.

\section{Temporal motifs}
\label{sec:motifs}

\paragraph{Static network motifs.}
The concept of network motifs was introduced by Milo et al.\ \citeyear{milo2002network} to describe small, recurring subgraph patterns that appear significantly more often in a real network than in a randomised reference. These patterns are considered the building blocks of network structure, so that different classes of networks (biological, social, technological) tend to be characterised by different sets of over-represented motifs \cite{milo2002network}. In static graphs, a motif is formally defined as a connected induced subgraph on $k$ nodes whose frequency in the observed network exceeds what would be expected by chance under a suitable null model \cite{milo2002network, newman2010networks}.


\paragraph{$\delta$-temporal motifs.}
Paranjape et al.\ \citeyear{paranjape2017motifs} extend the static motif concept to temporal networks by introducing $\delta$-temporal motifs.

\paragraph{The catalogue of $\delta$-temporal motifs.}

For the case of three nodes and three edges, Paranjape et al.\ 
\citeyear{paranjape2017motifs} enumerate a complete catalogue of 36 
$\delta$-temporal motifs, which together cover all distinct interaction patterns for three nodes. The 36 motifs 
contain three structural classes: four \emph{2-node motifs} (highlighted inside the green rectangle), in which 
all three edges occur between the same pair of nodes, 24 \emph{star motifs}, 
in which one center node interacts with two distinct neighbors and eight 
\emph{triangle motifs}, in which three nodes form a directed cycle 
\cite{paranjape2017motifs}. Within each structural class, the individual 
motifs are distinguished solely by the temporal ordering and direction of 
the three edges. The full catalogue is illustrated in 
Figure~\ref{fig:motifs_paranjape}.

\begin{figure}[h]
    \centering
    \includegraphics[width=\textwidth]{motifs_paranjape}
    \caption{All 36 two- and three-node, three-edge $\delta$-temporal motifs, 
    indexed as $M_{i,j}$ by row $i$ and column $j$. The first edge in each 
    motif originates from the green node and edge labels denote temporal ordering. 
    Figure from \citet{paranjape2017motifs}.}
    \label{fig:motifs_paranjape}
\end{figure}

\subsection{Counting algorithm}

Motif counting algorithms form the computational foundation of any motif-based analysis. This section first outlines the established methods for static graphs and then explains why a fundamentally different approach is required for temporal data.

\subsubsection{Static motif counting}

Static graph motif counting is a well-studied problem. Given a graph $G$, the goal is to count how many times a given subgraph pattern, such as a triangle or a directed cycle, appears in $G$. Algorithms for this task range from exact neighborhood-intersection methods to matrix-based approaches and sampling-based approximations, covering a broad spectrum of accuracy and efficiency trade-offs \cite{Yin2025}. These methods are computationally mature and underpin many downstream tasks such as clustering coefficient computation and community detection.

However, static methods operate on a time-aggregate representation: all edges that ever existed between two nodes collapse into a single connection, regardless of when they occurred \cite{kovanen2011temporal}. This means that the timely ordering is lost entirely. Paranjape et al \citet{paranjape2017motifs} show that even when a triangle is topologically present in a network, only a subset of its instances are temporally valid. The others either occur out of order or fall outside a given time window $\delta$ \cite{paranjape2017motifs}. A static counter treats all of them as identical. For blockchain transaction networks, this distinction is critical because coordinated malicious behavior happens through transaction sequences that have a specific temporal order, something that static aggregation cannot recover.

\subsubsection{Temporal motif counting}
Temporal motif counting extends the static counting problem by requiring that all edges of a motif occur in a specific order within a bounded time window $\delta$. 

The general counting algorithm proceeds in three steps: (1) identify all
static instances of the subgraph $H$ in the static graph $G$.
(2) for each instance, gather all corresponding temporal edges into a
time-sorted sequence $S$. and (3) process $S$ within a sliding window of
size $\delta$, tracking how often each ordered edge pattern appears
within that window \citep{paranjape2017motifs}. Algorithm~\ref{alg:paranjape}
formalizes this procedure.

\begin{figure}[t]
\begin{minipage}{\linewidth}

\begin{algorithm}[H]
\caption{$\delta$-Temporal Motif Counting adapted from \citet{paranjape2017motifs}}
\label{alg:paranjape}
\begin{algorithmic}[1]
\Require Time-sorted sequence $S = (e_1, \ldots, e_L)$,
         window $\delta$, motif length $l$
\Ensure Count of each $l$-edge $\delta$-temporal motif in $S$
\State $\textit{left} \gets 1$
\State $\textit{patternCounts} \gets \emptyset$
\For{$\textit{right} = 1, \ldots, L$}
    \While{$t_{\textit{right}} - t_{\textit{left}} > \delta$}
        \State \Call{RemoveFromWindow}{$e_{\textit{left}}$}
        \State $\textit{left} \gets \textit{left} + 1$
    \EndWhile
    \State \Call{AddToWindow}{$e_{\textit{right}}$}
\EndFor
\State \Return $\textit{patternCounts}$
\end{algorithmic}

\noindent\rule{\linewidth}{0.4pt}
\noindent\textbf{Procedure} \textsc{RemoveFromWindow}$(e)$
\hfill\textit{// expire $e$; withdraw prefix contributions}
\begin{algorithmic}[1]
\State $\textit{patternCounts}[e] \gets \textit{patternCounts}[e] - 1$
\ForAll{suffix $\sigma$ of length $l-1$ in $\textit{patternCounts}$}
    \State $\textit{patternCounts}[(e,\,\sigma)]
           \mathrel{-}= \textit{patternCounts}[\sigma]$
\EndFor
\end{algorithmic}

\noindent\rule{\linewidth}{0.4pt}
\noindent\textbf{Procedure} \textsc{AddToWindow}$(e)$
\hfill\textit{// admit $e$; extend all existing prefixes}
\begin{algorithmic}[1]
\ForAll{prefix $\pi$ of length $< l$ in $\textit{patternCounts}$,
        longest first}
    \State $\textit{patternCounts}[(\pi,\,e)]
           \mathrel{+}= \textit{patternCounts}[\pi]$
\EndFor
\State $\textit{patternCounts}[e] \gets \textit{patternCounts}[e] + 1$
\end{algorithmic}

\end{algorithm}
\end{minipage}
\end{figure}
 
For 2-node motifs this runs in 
$O(m)$ time, which is optimal. For 3-node, 3-edge star and triangle motifs, specialized algorithms reduce the worst-case complexity significantly, with the fast triangle algorithm achieving up to a 56.5$\times$ speedup over the general method on the Bitcoin transaction network \cite{paranjape2017motifs}. This scalability is relevant to the present work, as blockchain transaction graphs are large and the motifs we are interested in correspond directly to 3-node star and triangle patterns within a short time window.

\subsection{Motif fingerprint}
The identification of recurring subgraph patterns in a network provides insight into the local structural properties of individual nodes. However, raw motif counts aggregated over the entire graph provide only an overall description and do not make it possible to differentiate the behavior of individual nodes from one another. To address this, \citet{McDonnell2014} introduced the concept of a \textit{motif-role fingerprint}: a per-node vector $\mathbf{f}_v = [c_1, c_2, \ldots, c_k]$, where each component $c_i$ denotes the number of times node $v$ participates in a subgraph instance of motif type $M_i$. This distinguishes the fingerprint from a global motif count in that it retains which node participates and in which structural role. For instance, if a node acts as the initiator, intermediary, or receiver in a given pattern.

\citet{McDonnell2014} further show that functional motif-role fingerprints, obtained from counts of partial subgraphs, can be transformed into structural fingerprints by multiplying them with an invertible transformation matrix, establishing a mathematically well-defined correspondence between the two representations. The authors identify anomaly detection and the identification of structurally similar nodes as primary use cases, which makes it well applicable for this thesis.

The concept of a motif fingerprint as defined above applies to static directed graphs. In the context of this thesis, the notion is extended to temporal graphs following the $\delta$-temporal motif framework. The motif fingerprint of an address is therefore defined as:

\begin{equation}
    \mathbf{f}_v^{\delta} = \bigl[c_1^{\delta},\, c_2^{\delta},\, \ldots,\, c_k^{\delta}\bigr]
    \label{eq:motif_fingerprint}
\end{equation}

where $c_i^{\delta}$ is the number of times address $v$ participates in an instance of $\delta$-temporal motif type $M_i$. Practically, this is represented as a feature matrix indexed by Ethereum address, with one column per motif type. Each cell answers the question: how many times did this address participate in this type of time-local interaction pattern within the specified window $\delta$?

This serves two purposes. First, it reduces the high-dimensional transaction graph to a compact, interpretable, and fixed-size representation per address. Second, it encodes behavioral structure rather than raw transaction statistics because an address involved in coordinated activity is expected to exhibit a distinct fingerprint compared to ordinary users, as coordinated behavior by definition produces recurring, time-constrained interaction patterns among a group of addresses. The fingerprint therefore creates a behavioral signature of each address showing small recurring temporal subgraph patterns.

\section{Blockchain Relevance}
\label{sec:blockchain_relevance}

This section connects the foundations from the previous sections to the goal of this thesis, which is detecting coordinated malicious behaviour on Ethereum.

The abstractions introduced above are not only mathematically convenient, they also match the structure of Ethereum transaction data. As shown in the section on Ethereum fundamentals and in Section~\ref{sec:graph}, Ethereum transactions are directed interactions between addresses, and every interaction is recorded together with persistent metadata such as a timestamp, a value and the gas usage. The consequence for this thesis is that Ethereum is not just a static graph, but a transaction graph that keeps changing over time.

This temporal aspect is not a minor detail. Prior work on Ethereum anomaly detection argues that the time-related nature of transactions makes transaction graphs well suited for temporal network analysis \cite{pohlmann2019anomaly}. Studies of money trails report that activities such as phishing, gambling and money laundering show distinct cyclic transfer patterns that can only be identified when the order of transactions is preserved \cite{lal2021money}. Other work shows that suspicious behaviour often appears as rapid transaction bursts, sudden volume changes and short-window interaction patterns that a static graph cannot fully recover \cite{temporal_sampling_ethereum, ghoshcatalog2025}. Once the graph from Section~\ref{sec:graph} is extended with timestamps, it becomes possible to analyse not only \emph{who} interacts with whom, but also \emph{when}, \emph{how often} and \emph{in which order}.

As discussed in Section~\ref{sec:temporal_networks}, a static aggregation can show that a group of addresses is connected, but it cannot tell a slow, unrelated sequence of transfers apart from a tightly coordinated one that runs within minutes. This matters because coordinated malicious activity is rarely a single transaction. It usually consists of several transfers spread across multiple addresses, for example to hide the origin of funds, redistribute them or imitate legitimate activity.

For this thesis, this has a direct methodological consequence. The graph-theoretic and temporal foundations from the previous sections turn raw transaction histories into higher-level behavioural patterns. Temporal motifs identify recurring, time-local interaction structures, and motif fingerprints make these structures measurable for each address. The result is a representation that is well suited for detecting coordinated behaviour, because coordination is expected to show up not only in who is connected, but in repeated temporal interaction patterns among groups of addresses. Earlier network studies already model Ethereum as a directed transaction network of addresses and transfers \cite{wu2020temporal}. The step taken here is to treat that network as temporal.

This leads to the next step. If Ethereum is treated as a temporal transaction graph, the question is no longer only how addresses are connected, but which forms of coordinated malicious behaviour produce characteristic temporal interaction patterns. The following section therefore looks at coordinated malicious behaviour on Ethereum and describes the main categories considered in this thesis.

\subsection{Types of Coordinated Malicious Behaviour on Ethereum}
\label{sec:scam_types}
The Ethereum blockchain hosts a broad range of fraudulent and malicious activities, made possible by the combination of pseudo-anonymity, programmable smart contracts, and the irreversibility of on-chain transactions. A taxonomy of cybercrime on the Ethereum blockchain derived from 5,644 verified malicious addresses identifies 19 distinct categories of fraud, with the total stolen funds across the observation period reaching \$1.65 billion on-chain alone, a figure that exceeds self-reported losses to the U.S. Federal Trade Commission by a factor of almost 16, underscoring the scale of underreporting \citep{hornuf2025cybercrime}. This thesis focuses on the subset of malicious activity that is \emph{coordinated}, meaning it involves multiple addresses interacting in structured, time-local patterns. The following categories are most relevant.

\paragraph{Ponzi Schemes.}
Ponzi schemes are the largest category of on-chain cybercrime by aggregate stolen value, accounting for approximately 60\% of all funds transferred to verified malicious addresses on Ethereum \citep{hornuf2025cybercrime}. Operators promise high returns to investors and use incoming deposits from new participants to pay earlier ones, creating the appearance of a functioning investment operation \citep{bartoletti2020smart}. On Ethereum, Ponzi logic is often encoded directly in smart contracts, which execute the redistribution of funds automatically and transparently \citep{chen2020survey}. The transaction structure is inherently coordinated: funds flow from many victim addresses into a central contract address, which then redistributes them outward in a pattern that recurs as long as new participants enter. This creates characteristic fan-in and fan-out motif patterns in the transaction graph.

\begin{figure}[ht]
  \centering
  \begin{tikzpicture}
    \node[scamnode, fill=scamCoral,  label=above:{Investor 1}] (i1) at (-1.5,2.3) {};
    \node[scamnode, fill=scamPurple, label=above:{Investor 2}] (i2) at ( 1.5,2.3) {};
    \node[scamnode, fill=scamTeal,   label=below:{Contract}]   (c)  at ( 0,0)    {};
    % reciprocal pair between Investor 1 and Contract (deposit, then payout)
    \draw[scamedge] (i1) to[bend left=20] node[edgelbl,left]  {1} (c);
    \draw[scamedge] (c)  to[bend left=20] node[edgelbl,right] {3} (i1);
    % single deposit edge from Investor 2
    \draw[scamedge] (i2) to[bend right=8] node[edgelbl,right] {2} (c);
  \end{tikzpicture}
  \caption[Ponzi temporal motif]{Ponzi scheme. Two deposits arrive (edges~1,~2);
    the payout to the earlier investor (edge~3) fires only after a fresh deposit,
    so $t_1 < t_2 < t_3$.}
  \label{fig:motif-ponzi}
\end{figure}

\paragraph{Phishing.}
Phishing attacks involve deceiving users into approving malicious smart contract interactions or transferring funds to attacker-controlled addresses \citep{chen2020phishing}. In the Ethereum context, phishing has evolved beyond simple address impersonation to include signature-based token approval theft, where a victim signs a transaction that grants the attacker unlimited transfer rights over their tokens \citep{hornuf2025cybercrime}. The coordination aspect manifests in the use of drainer contracts: a central attacker-controlled contract collects funds from many victims in rapid succession, often within short time windows, producing a repeated fan-in temporal pattern anchored at the malicious contract address.

\begin{figure}[ht]
  \centering
  \begin{tikzpicture}
    \node[scamnode, fill=scamPurple, label=above:{Victim 1}] (v1) at (-2.2,2.2) {};
    \node[scamnode, fill=scamPurple, label=above:{Victim 2}] (v2) at (-0.6,2.2) {};
    \node[scamnode, fill=scamTeal,   label=below:{Drainer}]  (d)  at (-1.4,0)   {};
    \node[scamnode, fill=scamCoral,  label=below:{Attacker}] (a)  at ( 1.6,0)   {};
    \draw[scamedge] (v1) -- node[edgelbl,left]  {1} (d);
    \draw[scamedge] (v2) -- node[edgelbl,right] {2} (d);
    \draw[scamedge] (d)  -- node[edgelbl,above] {3} (a);
  \end{tikzpicture}
  \caption[Phishing temporal motif]{Phishing drainer. Victims are drained into a
    central contract within a tight window (edges~1,~2), which then sweeps the
    pooled funds to the attacker (edge~3).}
  \label{fig:motif-phishing}
\end{figure}

\paragraph{Pump-and-Dump Schemes.}
Pump-and-dump schemes involve a coordinated group of addresses accumulating a low-liquidity token, artificially inflating its price through simultaneous buy activity, and then selling their holdings once uninformed traders enter the market \citep{li2021cryptocurrency, hornuf2025cybercrime}. The coordination is explicit and time-critical: the buy phase and the sell phase each involve multiple addresses acting within narrow time windows in order to maximize price impact and exit before the market corrects. This produces dense, time-local clusters of transactions among a set of addresses, which are directly recoverable as temporal motifs.

\begin{figure}[ht]
  \centering
  \begin{tikzpicture}
    \node[scamnode, fill=scamCoral,  label=below:{Colluder 1}] (c1) at (-2.6,0) {};
    \node[scamnode, fill=scamTeal,   label=below:{Pool}]       (p)  at ( 0,0)   {};
    \node[scamnode, fill=scamCoral,  label=below:{Colluder 2}] (c2) at ( 2.6,0) {};
    \node[scamnode, fill=scamPurple, label=above:{Retail}]     (u)  at ( 0,2)   {};
    % reciprocal pairs: buy (1,2) then sell (4,5)
    \draw[scamedge] (c1) to[bend left=25]  node[edgelbl,above] {1} (p);
    \draw[scamedge] (p)  to[bend left=25]  node[edgelbl,below] {4} (c1);
    \draw[scamedge] (c2) to[bend right=25] node[edgelbl,above] {2} (p);
    \draw[scamedge] (p)  to[bend right=25] node[edgelbl,below] {5} (c2);
    % retail buys the pump, between the two phases
    \draw[scamedge] (u)  -- node[edgelbl,right] {3} (p);
  \end{tikzpicture}
  \caption[Pump-and-dump temporal motif]{Pump-and-dump. Each colluder buys
    (edges~1,~2) and later sells (edges~4,~5) against the same pool, with the
    retail victim's buy (edge~3) in between: $\{1,2\} \prec 3 \prec \{4,5\}$.}
  \label{fig:motif-pumpdump}
\end{figure}

\paragraph{Rug Pulls and Exit Scams.}
A rug pull occurs when the developers of a token project attract liquidity, typically by depositing token reserves into a decentralized exchange, and then abruptly withdraw all deposited funds, rendering the token worthless \citep{hornuf2025cybercrime, xia2021rugpull}. Exit scams follow a similar pattern at the level of initial coin offerings: funds raised from investors are transferred out of the project contract to attacker-controlled addresses immediately or within a short period after fundraising closes. Both patterns produce a concentrated outflow from a contract address to a small set of beneficiary addresses within a narrow time window, a temporal interaction structure that differs structurally from ordinary liquidity provision.

\begin{figure}[ht]
  \centering
  \begin{tikzpicture}
    \node[scamnode, fill=scamPurple, label=above:{LP 1}]      (l1)  at (-2.2,2.2) {};
    \node[scamnode, fill=scamPurple, label=above:{LP 2}]      (l2)  at (-0.6,2.2) {};
    \node[scamnode, fill=scamTeal,   label=below:{Pool}]      (p)   at (-1.4,0)   {};
    \node[scamnode, fill=scamCoral,  label=below:{Developer}] (dev) at ( 1.6,0)   {};
    \draw[scamedge]      (l1) -- node[edgelbl,left]  {1} (p);
    \draw[scamedge]      (l2) -- node[edgelbl,right] {2} (p);
    \draw[scamedgeheavy] (p)  -- node[edgelbl,above] {3} (dev);
  \end{tikzpicture}
  \caption[Rug pull temporal motif]{Rug pull / exit scam. Topologically identical
    to phishing, but in a larger time window. Liquidity is added gradually (edges~1,~2),
    then the whole pool is drained in one sudden transfer (edge~3).}
  \label{fig:motif-rugpull}
\end{figure}

\paragraph{Money Laundering.}
Money laundering on Ethereum typically involves moving illicitly obtained funds through a series of intermediate addresses to obscure their origin before depositing them at an exchange or converting them to a less traceable asset \citep{lal2021money}. This process, known as \emph{layering}, produces multi-hop transfer chains in which each hop occurs in rapid succession, as extended delay increases the risk of detection or address flagging. Temporal cycle analysis of Ethereum transactions has shown that money laundering exhibits distinct cyclic transfer patterns not observed in benign activity \citep{lal2021money}, making it a particularly strong candidate for temporal motif-based detection.

\begin{figure}[ht]
  \centering
  \begin{tikzpicture}
    \node[scamnode, fill=scamCoral, label=above:{Source}]   (s)  at (-3, 0.6) {};
    \node[scamnode, fill=scamTeal,  label=below:{Mule 1}]   (m1) at (-1,-0.6) {};
    \node[scamnode, fill=scamTeal,  label=above:{Mule 2}]   (m2) at ( 1, 0.6) {};
    \node[scamnode, fill=scamBlue,  label=below:{Exchange}] (e)  at ( 3,-0.6) {};
    \draw[scamedge] (s)  -- node[edgelbl,above] {1} (m1);
    \draw[scamedge] (m1) -- node[edgelbl,above] {2} (m2);
    \draw[scamedge] (m2) -- node[edgelbl,above] {3} (e);
  \end{tikzpicture}
  \caption[Money laundering temporal motif]{Money laundering (layering). Funds hop
    through intermediary nodes rapidly before cashing out (edges~1--3).
    The chain length~$k$ varies.}
  \label{fig:motif-laundering}
\end{figure}

\paragraph{Giveaway Scams.}
Giveaway scams involve the impersonation of reputable entities such as exchanges, public figures, or projects to solicit cryptocurrency transfers under the false promise of proportional returns \citep{hornuf2025cybercrime}. Structurally, they produce a rapid accumulation of inbound transfers from many victim addresses to a single attacker address within a short window, after which the address is typically abandoned. These constitute the second-largest category of Ethereum cybercrime by aggregate stolen value, with on-chain transfers totalling \$274 million in the period studied \citep{hornuf2025cybercrime}.

\begin{figure}[ht]
  \centering
  \begin{tikzpicture}
    \node[scamnode, fill=scamPurple, label=above:{Victim 1}] (v1) at (-2,2.3) {};
    \node[scamnode, fill=scamPurple, label=above:{Victim 2}] (v2) at ( 0,2.8) {};
    \node[scamnode, fill=scamPurple, label=above:{Victim 3}] (v3) at ( 2,2.3) {};
    \node[scamnode, fill=scamCoral]                          (a)  at ( 0,0)   {};
    \node[font=\small, below=1.5mm of a] (alab) {Attacker};
    \node[font=\footnotesize\itshape, text=gray, below=0.5mm of alab]
      {abandoned --- no outflow};
    \draw[scamedge] (v1) -- node[edgelbl,left]  {1} (a);
    \draw[scamedge] (v2) -- node[edgelbl,right] {2} (a);
    \draw[scamedge] (v3) -- node[edgelbl,right] {3} (a);
  \end{tikzpicture}
  \caption[Giveaway temporal motif]{Giveaway scam. A single inbound burst from
    many victims (edges~1--3); the diagnostic feature is the absence of any
    outgoing edge from the hub within the window.}
  \label{fig:motif-giveaway}
\end{figure}

\vspace{0.5em}
While the categories above differ in their mechanism and intent, they share a common structural property from the perspective of this thesis: all involve multiple addresses interacting in time-local, recurring patterns that leave distinguishable footprints in the temporal transaction graph. Table~\ref{tab:scam_types} summarizes the key characteristics of each category with respect to their temporal interaction structure.

\begin{table}[htbp]
\centering
\caption{Summary of coordinated malicious behaviour types on Ethereum and their characteristic temporal interaction patterns.}
\label{tab:scam_types}
\begin{tabular}{p{3.2cm} p{5.2cm} p{4.8cm}}
\toprule
\textbf{Type} & \textbf{Mechanism} & \textbf{Temporal Pattern} \\
\midrule
Ponzi Scheme     & Redistribution of new deposits to old investors via contract & Cyclic fan-in / fan-out at contract address \\
Phishing         & Approval theft or direct transfer to drainer contract         & Rapid fan-in to single attacker address \\
Pump-and-Dump    & Coordinated buy/sell in narrow time window                    & Dense bidirectional cluster among colluding addresses \\
Rug Pull         & Liquidity withdrawal from exchange contract                   & Concentrated outflow from contract to beneficiary set \\
Money Laundering & Multi-hop layering to obscure fund origin                     & Sequential chained transfers with short inter-hop delays \\
Giveaway Scam    & Impersonation to solicit inbound transfers                    & Short-window fan-in followed by address abandonment \\
\bottomrule
\end{tabular}
\end{table}


\section{Summary}

- 5-8 pages




\chapter{Method} 
\label{ch:method}

\section{System overview}
\label{sec:system_overview}

This chapter describes how the detection approach is put into practice. The
method is implemented as a single program that turns raw Ethereum transaction
records into a decision about whether each address behaves maliciously. The
program is not used interactively. It runs a fixed sequence of steps, always in
the same order, so that the same data and the same settings always lead to the
same result.

The program runs in two phases. A \emph{preparation} phase reads the raw data
once, decides which addresses take part in the study, and fixes how they are
divided into a development set and a test set. An \emph{execution} phase then carries
out the actual detection on top of these prepared data. Keeping the two apart
means that the data and the evaluation setup are settled before any model is
trained, so the reported results cannot be shaped, even by accident, by later
changes to the data or the settings. The execution phase is checkpointed as
well. After an expensive motif count has been committed successfully, an
interrupted run verifies and reuses it instead of starting that count again.

Overall the program moves from raw transactions to results in eight steps, shown
in Figure~\ref{fig:method_steps}. The first two belong to the preparation phase,
the remaining six to the execution phase.

\begin{enumerate}
  \item \textbf{Loading the data.} The program reads the transaction records and
        the address labels of the chosen dataset. For the main dataset the
        transaction table is very large, with several hundred million transfers,
        so it is read in small batches instead of being held in memory all at
        once. The datasets are described in Section~\ref{sec:data_acquisition}.

  \item \textbf{Preparing the data.} Not every labelled address can be used. The
        program measures how active each address is and keeps only those with
        enough transactions to take part in a motif. It also sets aside a small
        number of extremely active hub addresses, so that a few exchange-sized
        accounts cannot dominate the analysis. The remaining addresses are split
        into a development set, used for tuning, and a separate test set, used
        only for the final measurement. For each address the program stores just
        the transactions that involve it and writes these prepared tables to
        disk, so that every later step works on exactly the same data.

  \item \textbf{Building the temporal graph.} The prepared transactions are
        loaded into a temporal graph, in which every address is a node and every
        transfer is a directed edge that carries the time at which it happened.
        Repeated transfers between the same two addresses are kept separately,
        because their number and their order are what the motifs are built from.
        This step is described in Section~\ref{sec:graph_construction}.

  \item \textbf{Counting the motifs.} The program first runs exact synthetic
        fixtures as a correctness gate. It then counts for each address how
        often each temporal motif occurs among the transactions around it.
        Because a motif only counts when its transactions fall close enough
        together in time, the counting is repeated for a range of time windows,
        from five minutes up to one week. This step is described in
        Section~\ref{sec:motif_extraction}. Each completed count is stored as an
        immutable checkpoint before the next one begins.

  \item \textbf{Building the fingerprints.} The motif counts of an address form
        its fingerprint, a short list of numbers that records how often the
        address took part in each motif. From the same transactions the program
        also computes a set of simpler behavioural features, such as how much
        value the address sent and received, how many partners it had, and how
        its activity was spread over time. This produces three views of every
        address: the simpler features on their own, the motif fingerprint on its
        own, and the two together.

  \item \textbf{Running the detection.} For each of the three views the program
        trains a classifier to tell malicious and benign addresses apart. It
        uses the development set to choose the time window that works best and
        then measures how well the classifier performs on the untouched test
        set. For RQ3 the same idea is used, except that the classifier sorts
        criminal addresses into their behaviour category rather than making a
        malicious-or-benign decision. The features and models are described in
        Section~\ref{sec:detection_methods}.

  \item \textbf{Analysing the results.} The program then answers the research
        questions. It compares the three views to see whether the motif
        fingerprint improves detection over the simpler features (RQ1) and which
        time window makes coordinated behaviour most visible (RQ2), and it checks
        whether different categories of malicious behaviour have distinct
        fingerprints and their own best time windows (RQ3). Each comparison is
        supported by statistical tests, so that a difference is not mistaken for
        chance, and the whole run is repeated under several orderings of
        transactions that share the same timestamp, to confirm that the outcome
        is stable. The evaluation is described in
        Section~\ref{sec:evaluation_protocol}.

  \item \textbf{Visualising the results.} Finally the program writes the tables
        and figures used in the Results chapter, including the comparison of the
        three views, the performance curve across the different time windows, and,
        for RQ3, how the motif fingerprints and best windows differ between
        behaviour categories.
\end{enumerate}

\begin{figure}[ht]
  \centering
  \ResGraphic{0.75}{05-method-steps}
  \caption{The eight steps of the pipeline. The first two make up the preparation
  phase, and the remaining six the execution phase, shown here as its
  representation, detection, and analysis stages.}
  \label{fig:method_steps}
\end{figure}

The first two research questions run on the main binary dataset and the third on
the category-labelled dataset. The following sections describe the most important
of these steps in more detail, starting with the data.

\section{Data acquisition}
\label{sec:data_acquisition}

Studying coordinated malicious behaviour on-chain needs two things at once:
transaction data, and reliable labels that say which addresses are malicious.
This combination is scarce. Transactions are public, but the intent behind an
address can only be established off-chain, through abuse reports, exchange
records, or investigations, and such labels are rare and unevenly spread across
the network \citep{shi2025realcats}. To work around this, the thesis draws on
four datasets, each in a different role. One main dataset carries the
evaluation, a second Ethereum dataset serves as an independent check, a
real-world dataset adds fine-grained scam categories, and a small set of
synthetic graphs is used only to confirm that the motif counting is correct. The
origin and prior use of the two published datasets, BLTE and Real-CATS, were
already covered in Section~\ref{sec:relwork_datasets}. Here each dataset is
described in the same way, by its source, its contents, and its suitability for
this thesis.

\paragraph{HuggingFace Ethereum Fraud Dataset.}
\label{sec:dataset_hf}
\textit{Source.} The dataset is published on the Hugging Face Hub
\citep{fesevu2025ethfraud}. This thesis uses its full-history release, which
pairs a complete transaction table with a set of labelled addresses, rather than
one of the smaller monthly extracts.

\textit{Contents.} The dataset has two parts. The first is a transaction table
of 380{,}566{,}661 transfers recorded between August 2015 and August 2025. Each
transfer carries:
\begin{itemize}
  \item the sending and receiving address, as hashed identifiers,
  \item the timestamp of the transfer,
  \item the transferred value,
  \item the transaction hash and block number,
  \item a flag marking contract-creation transactions, and
  \item the transaction fee.
\end{itemize}
The second part is a label table of 114{,}739 addresses, of which 60{,}086 are
benign and 54{,}653 fraudulent. Each labelled address also states whether it is
a contract or an ordinary account, and carries a short text description of the
reported activity.

\textit{Suitability.} This is the main dataset of the thesis, and it carries the
evaluation for RQ1 and RQ2. It is the only dataset used here that
combines full timestamped transaction histories with a large and roughly
balanced set of malicious and benign labels. The timestamps matter because
temporal motifs depend on the exact timing and order of transfers, and the size
of the label set is large enough to draw more than one cohort of addresses while
still keeping an untouched test set aside.

\paragraph{BLTE, Benchmark Labelled Transactions of Ethereum.}
\label{sec:dataset_blte}
\textit{Source.} BLTE was introduced by \citet{alemari2021labeled} and is built
on the Ethereum Classic chain. Its origin and the results reported on it by
earlier work were described in Section~\ref{sec:relwork_datasets}. Here it is
used as obtained from the public release.

\textit{Contents.} The dataset contains 71{,}250 transactions from October 2017
to December 2019, labelled one transaction at a time as normal or abnormal.
There are 57{,}000 normal and 14{,}250 abnormal transactions, so about one
transaction in five is abnormal. Each record provides:
\begin{itemize}
  \item the transaction hash and nonce,
  \item the sending and receiving address,
  \item the transferred value, the gas, and the gas price,
  \item the block timestamp, block number, and block hash, and
  \item four labelling fields that mark whether the sender or the receiver is a
        known scam, and of which category.
\end{itemize}
A set of pre-computed per-address features is distributed alongside the raw
transactions, but this thesis does not use them, for the reason given below.

\textit{Suitability.} BLTE is used as a second, independent Ethereum dataset, to
check whether the findings on the main dataset also appear elsewhere. Its labels
are very imbalanced at the address level: of 73{,}034 active addresses, only 169
are marked as scams. It is therefore treated as a secondary check rather than a
headline result, and the imbalance is handled by weighting the classes during
training. The pre-computed features are deliberately left aside, so that every
dataset is described by the same behavioural features, computed in the same way,
which keeps the comparison with the main dataset fair.
BLTE is prepared and executed as a separate optional run. Its inclusion or
exclusion does not determine whether the primary HF and Real-CATS runs are
citable.

\paragraph{Real-CATS, Real-world Criminal Account Transaction Set.}
\label{sec:dataset_realcats}
\textit{Source.} Real-CATS was introduced by \citet{shi2025realcats}. Its design
and the sources of its labels were described in
Section~\ref{sec:relwork_datasets}. This thesis uses its Ethereum portion
together with the detailed transaction records released alongside it.

\textit{Contents.} The Ethereum portion covers 12{,}572 criminal and 16{,}019
benign addresses. For each address the dataset provides its transaction history,
split into five separate layers: ordinary transfers, internal transfers, and
ERC-20, ERC-721, and ERC-1155 token transfers. What sets Real-CATS apart from
the other datasets is that its criminal addresses carry a specific category,
rather than a single malicious flag. The categories are:
\begin{itemize}
  \item Phishing Scam and Impersonation Scam,
  \item Contract Exploit Scam and Hack Scam,
  \item Blackmail Scam and Fake Returns Scam,
  \item Rug Pull Scam, NFT Airdrop Scam, and Pigbutchering Scam.
\end{itemize}

\textit{Suitability.} Real-CATS is the dataset used for RQ3. Its detailed
categories are what make that question possible, because they allow the motif
fingerprints of different kinds of coordinated behaviour to be compared directly,
rather than only a single malicious-versus-benign split. The labels are also
independent of the main dataset, since they come from public abuse reports
instead of the process behind the HuggingFace release. The categories are very
unevenly filled, with phishing and impersonation together accounting for most of
the criminal addresses, so RQ3 is limited to the categories that have enough
addresses to support a fair comparison and treats the smaller ones as
descriptive only.

\paragraph{Synthetic fixtures.}
\label{sec:dataset_synthetic}
\textit{Source.} The synthetic data is generated in code. Each fixture is a
small graph of exactly three transfers, built by hand to represent one known
motif.

\textit{Contents.} The fixtures contain no real transactions. Together they
cover all 40 positions in Raphtory's returned catalogue, including the eight
triangle roles that are checked and then excluded from the target-incident
experiments. They also cover the boundary cases where transactions fall exactly
on the edge of a time window or share the same timestamp.

\textit{Suitability.} These graphs are used only to confirm that the motif
counting is correct, meaning that a known pattern is recognised as the motif it
is meant to be. They are never used to measure detection performance, and are
not compared against the real datasets.

\begin{table}[ht]
\centering
\caption{The datasets used in this thesis and their roles. The malicious share
is computed over labelled addresses.}
\label{tab:dataset_comparison}
\footnotesize
\setlength{\tabcolsep}{4pt}
\begin{tabular}{@{}p{2.7cm} p{2.5cm} p{2.4cm} p{2.6cm} p{1.9cm}@{}}
\toprule
\textbf{Property} & \textbf{HuggingFace} & \textbf{BLTE} & \textbf{Real-CATS} & \textbf{Synthetic} \\
\midrule
Chain              & Ethereum      & Ethereum Classic & Ethereum      & generated \\
Transactions       & 380.6\,M      & 71{,}250         & per-address histories & 3 per fixture \\
Label level        & address       & transaction      & address       & motif \\
Labelled addresses & 114{,}739     & 73{,}034 active  & 28{,}591      & n/a \\
Malicious share    & 47.6\,\%      & 0.23\,\%         & 44.0\,\%      & n/a \\
Date range         & 2015--2025    & 2017--2019       & per release   & controlled \\
Category labels    & no            & no               & yes (9)       & no \\
Role               & RQ1, RQ2      & RQ1 (secondary)  & RQ3           & validation \\
\bottomrule
\end{tabular}
\end{table}

\section{Temporal graph construction}
\label{sec:graph_construction}

The prepared transactions are turned into a temporal graph, which is the object
the motifs are later counted on. This section describes the graph model, how the
graph is actually built from each dataset, and the simple features that are read
straight from it and later serve as a comparison point for the motifs.

\subsection{Graph model}

Each address is represented as a node and each transfer as a directed edge that
carries the time at which it happened. The result is a temporal directed
multigraph
\[
  \mathcal{G}=(V,E,\tau,w),
\]
where $V$ is the set of addresses, $E$ the set of transfers, $\tau(e)$ the
timestamp of transfer $e$, and $w(e)$ its transferred value. Two properties are
kept on purpose. Parallel transfers between the same pair of addresses are not
merged, because how many there are and in which order they occur is exactly what
the motifs are built from. Edge direction is kept for the same reason, since an
address that only receives, one that only sends, and one that does both leave
different motif patterns behind.

\subsection{Building the graph from the data}

HuggingFace and Real-CATS are stored in very different shapes, so the graph is
assembled differently for each, but always without holding more data in memory
than needed.

For HuggingFace the transaction table is far too large to load as a whole, with
more than 380 million transfers. It is therefore read in small batches, first to
measure how active each labelled address is and then to keep only the transfers
that touch one of the selected addresses. The graph used for the experiment is
this collection of transfers around the target addresses, rather than the entire
network.

This choice has one consequence worth stating. Because every kept transfer
touches a target address, the graph captures every motif that is centred on a
target, as well as every back-and-forth between a target and a single partner.
It does not reliably capture triangles, which are motifs among three addresses
where one edge runs between two of the target's partners and so need not touch
the target itself. Triangle motifs are therefore left out, which leaves the 32
motif roles described in Section~\ref{sec:motif_extraction}.

Real-CATS is stored as one transaction history per address, and the same
transfer can appear in the histories of several addresses. These repeated copies
are removed once, using a disk-backed index so that the de-duplication does not
depend on everything fitting in memory. Each history is also split into five
layers, namely ordinary transfers, internal transfers, and ERC-20, ERC-721 and
ERC-1155 token transfers. Every layer becomes its own set of edges, and motifs
are counted within a layer and never across layers, since two transfers recorded
in different layers do not form a genuine time-ordered sequence.

\subsection{The non-motif baseline}

RQ1 asks whether motifs add anything beyond ordinary address behaviour, so it
needs a description of each address that does not use motifs at all. This
baseline is computed from the same transfers and the same addresses as the motif
fingerprints. Table~\ref{tab:results_baseline_features} lists the resulting 17
features, grouped into how active an address is, how much value it moves, the
structure of its immediate neighbourhood, and how its activity is spread over
time.

\begin{table}[ht]
\centering
\caption{The non-motif baseline features used for RQ1.}
\label{tab:results_baseline_features}
\small
\begin{tabular}{@{}p{3.4cm} p{9.5cm}@{}}
\toprule
\textbf{Feature group} & \textbf{Features} \\
\midrule
Activity &
Number of transfers sent, received, and in total \\[2pt]
Transferred value &
Outgoing and incoming sum, mean, and sample standard deviation, plus the total
value moved \\[2pt]
Neighbourhood structure &
Number of distinct outgoing neighbours, number of distinct incoming neighbours,
and reciprocity, the share of an address's neighbours that it both sends to and
receives from \\[2pt]
Relative timing &
Active span and the mean, median, and sample standard deviation of the gaps
between successive transfers that touch the address \\
\bottomrule
\end{tabular}
\end{table}

A few choices keep this baseline honest. Absolute first and last transaction
times are deliberately left out, because they mostly reveal when the data was
collected or how old a campaign is rather than how an address behaves, and a
model could otherwise separate the classes by collection window instead of by
behaviour. Measures that describe an address's position in the wider network are also left
out, because computing them reliably would need more of the surrounding graph
than is kept here. Where a value cannot be computed, for instance a measure of
variation for an address with too few transactions, it is set to zero.

\section{Temporal motif extraction}
\label{sec:motif_extraction}

The baseline describes an address by totals and averages, which cannot see the
order in which its transfers happen. Two addresses can have the same number of
transfers, the same partners, and the same active span, yet behave very
differently: one might receive money from many senders within minutes and pass
it straight on, while the other collects the same money slowly over months.
Motifs are what capture this difference. Recall from Section~\ref{sec:motifs}
that a $\delta$-temporal motif is a short, time-ordered sequence of transfers
that all fall within a window of length $\delta$, and that stacking an address's
motif counts into a vector gives its motif fingerprint
(Equation~\ref{eq:motif_fingerprint}). This section describes how those
fingerprints are counted and turned into features.

\subsection{Motif counting and retained roles}

The counting itself is done by Raphtory, whose
\texttt{local\_temporal\_three\_node\_motifs} routine is written in Rust and
processes addresses in parallel. For every address it returns a \emph{local}
count for each motif role, that is, the number of $\delta$-temporal instances in
which the address plays that particular part. A star motif is counted only for
the address at its centre, and a two-node motif for both of the addresses
involved.

Raphtory reports 40 such roles per address: 24 star roles (three classes of
eight in-and-out direction patterns each), 8 two-node roles, and 8 triangle
roles. These are local roles, so the same global two-node motif is split by which
endpoint and direction is meant, which is why there are 40 local roles rather
than the 36 global motifs shown in Figure~\ref{fig:motifs_paranjape}. This work
keeps the 24 star and 8 two-node roles, which gives a 32-dimensional fingerprint
per address, and drops the 8 triangle roles for the reason given in
Section~\ref{sec:graph_construction}: the graph only keeps transfers that touch a
target address, so it can miss the third edge of a triangle, the one that runs
between two of the target's partners. Counting triangles would then produce zeros
that reflect this missing edge rather than the address's behaviour. Raphtory has no
option to skip motif classes, so all 40 are computed and the eight triangle
columns are removed afterwards. An address that takes part in no counted motif
keeps a fingerprint of all zeros rather than being dropped from the study.

The retained roles fall into a few structural classes, each with a
straightforward reading in a transaction network.
Table~\ref{tab:motif_classes} lists them together with the kind of coordinated
behaviour they tend to reflect.

\begin{table}[h]
\centering
\caption{Retained structural motif classes (after \cite{paranjape2017motifs})
and their behavioural interpretation in Ethereum transaction networks. Each
class covers several $\delta$-temporal roles that differ in the direction and
order of the three transfers.}
\label{tab:motif_classes}
\begin{tabular}{p{3.5cm} p{4.5cm} p{5cm}}
\toprule
\textbf{Class} & \textbf{Structure} & \textbf{Fraud interpretation} \\
\midrule
Two-node, same direction &
    $A \to B,\; A \to B,\; A \to B$ &
    Repeated transfers, which may indicate instalment-based extraction \\[4pt]
Two-node, mixed direction &
    $A \to B,\; B \to A,\; A \to B$ &
    Rapid back-and-forth, as in wash-trading or fee manipulation \\[4pt]
Star, all incoming &
    $B \to A,\; C \to A,\; C \to A$ &
    Funds consolidating at a single hub, as in money-mule collection \\[4pt]
Star, all outgoing &
    $A \to B,\; A \to C,\; A \to C$ &
    Mass distribution from one source, as in airdrop farming or Ponzi payouts \\[4pt]
Star, mixed &
    $A \to B,\; C \to A,\; A \to C$ &
    A relay node forwarding funds, as in a layering intermediary \\
\bottomrule
\end{tabular}
\end{table}

\subsection{Ordering transfers with equal timestamps}

Because motifs are ordered sequences, the transfers have to be put in a definite
order before they can be counted. Time supplies most of that order, but many
transfers carry the same timestamp because they sit in the same block, and time
alone cannot separate them. When the data records where a transfer sits within its
block, that finer order is used. The HuggingFace data does not record this, so its
transfers that share a timestamp have no fixed order. To make sure the results do
not rest on an arbitrary choice, the whole analysis is run under seven versions of
the order. All of them sort by time first, so the time order of the motifs is
never changed, and they differ only in how they arrange the transfers that share a
timestamp: one fixed order, the same order reversed, and five random shuffles.
Section~\ref{sec:evaluation_protocol} reports whether the conclusions hold under
all seven. The order is computed without loading the whole table into memory, and
the sorted transfers are passed to Raphtory in large batches.

\subsection{Fingerprint representations}

The retained counts are stored under stable names such as \texttt{star\_pre\_iii}
and \texttt{two\_node\_ioi}. For RQ1 and RQ2 each count $c$ is replaced by
$\log(1+c)$. A few very active addresses reach counts in the thousands while most
stay in the low single digits, and on the raw scale these few large numbers would
dominate the model. Taking the logarithm compresses this range, and adding one
first means that a count of zero stays zero. These 32 transformed values are the
motif fingerprint used by the detection step, either on their own or joined to the
baseline of Section~\ref{sec:graph_construction}. How the fingerprint and the baseline are
combined into the feature sets compared in RQ1 is described in
Section~\ref{sec:detection_methods}.

RQ3 keeps the five Real-CATS layers separate. Because the layers differ greatly
in how many transfers they contain, comparing raw counts would let the busiest
layer dominate. Instead, within each layer the counts are turned into proportions by dividing each
count by the layer total, and the square root of each proportion is taken. This is
known as a Hellinger transform:
\[
    h_{vml} = \sqrt{\frac{c_{vml}}{\sum_j c_{vjl}}},
\]
where $c_{vml}$ is the count of role $m$ for address $v$ in layer $l$. Using
proportions removes the effect of volume, so only the mix of motifs remains, and
the square root keeps a few large proportions from dominating, much as the
logarithm does for the counts. A layer with no transfers is kept as all zeros. The
same counts in log form, as used for RQ1 and RQ2, are also kept as a check.

For every window the pipeline additionally records whether each address has any
nonzero count at all. This coverage measure separates a real improvement from a
more informative time scale from the trivial effect of fingerprints simply
becoming nonzero as the window grows.

\subsection{Temporal-window sweep}

The counting is repeated over a fixed grid of seven windows,
\[
    \delta \in \{300,\,1800,\,3600,\,21600,\,86400,\,259200,\,604800\}
    \text{ seconds},
\]
which correspond to five minutes, 30 minutes, one hour, six hours, one day, three
days, and one week. All windows use the same addresses and the same train and
test split, and the window is chosen only on the development data, as described
in Section~\ref{sec:evaluation_protocol}. RQ2 studies how detection changes
across this grid, and RQ3 additionally asks whether different categories of
behaviour are best seen at different windows.

\section{Detection methods}
\label{sec:detection_methods}

The previous two sections produced two descriptions of every address: the
non-motif baseline of Section~\ref{sec:graph_construction} and the motif
fingerprint of Section~\ref{sec:motif_extraction}. This section describes how
those descriptions are turned into a detector and which models are used. The
binary detection for RQ1 and RQ2 is kept apart from the category detection for
RQ3, because the two ask different questions and use different datasets.

\subsection{The three feature views}

RQ1 asks whether the motif fingerprint adds anything beyond ordinary address
behaviour. The cleanest way to answer this is a controlled comparison, in which
only the features change and everything else stays fixed. Three views of each
address are therefore built from the two descriptions above:
\begin{itemize}
  \item the \emph{baseline} view, the 17 non-motif features of
        Table~\ref{tab:results_baseline_features},
  \item the \emph{motifs} view, the 32 log-transformed motif counts, and
  \item the \emph{combined} view, the two simply placed side by side.
\end{itemize}
All three views use exactly the same addresses and the same train and test split,
and no address is dropped just because its fingerprint is all zeros. Because the data is
held constant, any difference in detection between the views can only come from
the features, which is what makes this a fair test of RQ1.

RQ2 then asks at which time window the motifs are most telling. For this the
combined view is evaluated at each of the seven windows from
Section~\ref{sec:motif_extraction}. The baseline does not depend on the window, so
it stays the same throughout and serves as a fixed reference line.

\subsection{Classifiers for RQ1 and RQ2}

The detector for the binary questions is a random forest, an ensemble of decision
trees that each vote on whether an address is malicious. A random forest suits
this task because it can pick up interactions between different motif roles
without any manual feature scaling, and because it is a strong, well-understood
default rather than a model tuned hard for the highest possible score. Its
settings are fixed in advance, with 100 trees, a limited depth, balanced class
weights, and a fixed random seed, and they are not adjusted on the test data. The
forest is trained on the development addresses and then scores each address in the
test set with a probability of being malicious. The time window is chosen on the
development data by cross-validation, using the rule described in
Section~\ref{sec:evaluation_protocol}.

As a robustness check, a gradient-boosted tree model (XGBoost) is also fitted, on
the combined view at the chosen window. If it agrees with the random forest, that
strengthens confidence in the result, but the RQ1 decision always rests on the
random forest, which is fixed beforehand. Table~\ref{tab:results_models}
summarises all three models.

\begin{table}[ht]
\centering
\caption{The three models used in the detection experiments.}
\label{tab:results_models}
\small
\begin{tabular}{@{}p{2.8cm} p{3.7cm} p{6.6cm}@{}}
\toprule
\textbf{Task} & \textbf{Model} & \textbf{Fixed settings} \\
\midrule
RQ1 and RQ2 (main) &
Random forest &
100 trees, depth 10, balanced class weights, a fixed seed, allocated CPU threads \\[2pt]
RQ1 (robustness) &
Gradient-boosted trees (XGBoost) &
100 rounds, depth 6, learning rate 0.1, positive-class weight from the training
split, a fixed seed \\[2pt]
RQ3 &
Multinomial logistic regression &
Standardised features, balanced class weights, up to 2{,}000 iterations, a fixed
seed \\
\bottomrule
\end{tabular}
\end{table}

\subsection{Classifier for RQ3}

RQ3 does not separate malicious from benign addresses. It sorts criminal
addresses into their behaviour category, so the model has one class per category
rather than two. Only categories with at least 100 addresses take part, because the categories are
very unevenly filled, as explained in Section~\ref{sec:data_acquisition}, and
categories are never merged into larger groups.

The input here is the per-layer fingerprint of
Section~\ref{sec:motif_extraction}, expressed as motif proportions, which captures
the mix of motifs an address shows in each transaction layer. The model is a
logistic regression in its multi-class form. Its inputs are first put on a common
scale, and this scaling is learned only from the training data so that the test
data cannot influence it. A logistic regression is chosen deliberately, because it
is a simple model with few parameters. This tests whether the motif proportions
alone can separate the categories, without a more complex model that might overfit
the smaller categories.

To answer both halves of RQ3, two kinds of window are selected on the development
data. A single global window is chosen to separate all categories as well as
possible at once, and a separate window is chosen for each category to see it at
its best. The final model is refitted on all development addresses and gives each address in
the test set one probability per category.

\subsection{Handling class imbalance}

The classes are uneven to different degrees across the experiments, only slightly
for the main dataset and strongly for BLTE and for the Real-CATS categories. This
is handled in a single way, by weighting each class in inverse proportion to how
often it appears, so that a rare class counts for as much as a common one during
training. The random forest and the logistic regression use balanced class
weights, and the gradient-boosted model uses the matching positive-class weight
from its training split. No artificial examples are added, and the number of
addresses in each class is left unchanged.

\section{Evaluation}
\label{sec:evaluation_protocol}

\subsection{Scope}

The evaluation measures how well an address can be classified after the fact, from
its full recorded history. It does not estimate how early an address could have
been caught, and it makes no causal claim about how a fraud was carried out. As
described in Section~\ref{sec:system_overview}, the data and every setting are
fixed in a preparation step before any model is trained, so nothing in the
evaluation can be tuned on the test data. The parts that enforce this, such as the
recorded settings and the checks that the code and data have not changed, are
described in Section~\ref{sec:implementation}.

\subsection{Cohorts and splits}

RQ1 and RQ2 draw two separate groups of addresses from the HuggingFace dataset. An
address is eligible only if it takes part in at least three transactions, since
fewer cannot form one of the three-edge motifs. Before the groups are drawn, the
most extreme addresses are removed: within each combination of label, contract
type, and activity level, the top 0.5\,\% most active addresses are dropped.
Without this, a few exchange-scale hubs would dominate the graph on their own. The
cut is made within each combination rather than across the whole dataset, because
the busiest addresses are almost all benign exchange wallets, so a single global
cut would remove far more benign than fraudulent addresses and quietly shift the
balance between the classes. The rule is fixed before any model is trained, and
every removed address is recorded.

The two groups serve different purposes. The \emph{prevalence} cohort is a plain
random sample of up to 5{,}000 eligible addresses, which keeps the natural mix of
fraudulent and benign addresses and checks that the findings hold on a realistic
distribution. The \emph{controlled} cohort is drawn from the remaining addresses,
up to 250 per combination of label, contract type, and activity level, which
balances the classes and holds activity and contract type roughly constant so that
a difference cannot be explained by those alone. The two cohorts share no
addresses and are evaluated as separate experiments.

Splitting the addresses needs one precaution. Some fraud arrives in campaigns of
many near-identical addresses, and if some of them landed in the training set
while their near-twins landed in the test set, the detector would look better than
it really is. Known campaign families are therefore kept together as one group,
and every split keeps a whole group on a single side. Each cohort is divided about
80 to 20 into a development set and a test set with a group-aware, stratified
split, and the development set is divided again into repeated five-fold splits,
under different random seeds, for choosing the time window. Cohort membership is
included in the stratification, while campaign assignments are made across both
cohorts together; a campaign represented in both cohorts therefore remains on
the same side and in the same validation fold. RQ3 uses its own 80 to
20 split, stratified by category, again keeping only categories with at least 100
addresses and never merging them.

\subsection{Metrics and window selection}

The main measure is average precision, which summarises how well the detector
ranks malicious addresses above benign ones across all decision thresholds. It
suits imbalanced data because it focuses on the rarer positive class. F1,
precision, and recall are reported as secondary measures at a fixed threshold of
0.5, which is never tuned on the test set, and ROC-AUC is reported alongside them.
For RQ3 the matching measures are the macro-averaged and per-category average
precision, the macro F1, and a confusion matrix.

The time window for RQ1 is chosen on the development data alone. Each window is
scored by cross-validation, and the choice follows the one-standard-error rule:
take the window with the best average score, then pick the shortest window whose
score is still within one standard error of it. This prefers the shortest window
that scores as well as the best within the margin of error, rather than the single
highest-scoring window, whose lead may be noise. The test set plays no part in this
choice.

\subsection{Statistical tests}

For RQ1 the effect is the difference in test-set average precision between the
combined view and the baseline. Two tests support it. The first is a permutation
test: it repeatedly swaps whole campaigns between the two models, recomputes the
difference, and does this 10{,}000 times to build up the range of differences that
could appear by chance. This gives a one-sided $p$-value for the claim that the
combined view does not improve on the baseline. The second is a bootstrap: it
re-draws the campaigns at random 10{,}000 times to give a 95\,\% confidence
interval for the difference. Both work on whole campaigns rather than single
addresses, so that addresses from the same campaign always stay together.

For RQ2 the chosen window is compared with each of the other windows using the
same kind of bootstrap. Because many windows are compared at once, the resulting
$p$-values are adjusted for multiple testing with the Benjamini--Hochberg
procedure, and a window is declared the single best only if it still beats every
alternative afterwards.

For RQ3 two things are tested. The first is whether the categories can be separated
at all. The true category labels are shuffled at random 10{,}000 times, and if the
model does no better on the real labels than on the shuffled ones, the fingerprints
carry no signal about the category. The second is whether a category is seen better
at its own best window than at the shared one. This is checked for each category
with a bootstrap, and because it is repeated for every category, the results are
again adjusted for multiple testing.

\subsection{Robustness to event ordering}

Every result is produced under each of the seven orderings of
Section~\ref{sec:motif_extraction}. An ordering-dependent conclusion, such as an
RQ1 improvement or an RQ2 best window, is accepted only when its direction and its
decision agree across all seven orderings. This shows that a conclusion does not
depend on how transfers sharing a timestamp are ordered. It does not recover the
real order of those transfers, which the data does not record in the first place.

\section{Implementation details}
\label{sec:implementation}

This section records the tools the pipeline is built on, how the code is
organised, the main engineering decisions behind it, and the measures that make a
run reproducible.

\subsection{Language and libraries}

The pipeline is written in Python, with the temporal motif counting carried out by
Raphtory's Rust extension. Everything runs on the CPU, so no graphics card is
needed, and Raphtory, DuckDB, the random forest, and XGBoost are all set up to use
the thread allocation recorded for the run. Native thread limits are set before
the numerical libraries are imported, so two coordinated workers cannot each
claim every core. Development and validation used
Python~3.13 in an isolated virtual environment. The exact interpreter and package
versions of each run are recorded in its protocol, and that record, rather than the
environment described here, is what counts for the final numbers.
Table~\ref{tab:libraries} lists the main libraries and their role.

\begin{table}[h]
\centering
\caption{The main libraries used by the pipeline.}
\label{tab:libraries}
\renewcommand{\arraystretch}{1.3}
\small
\begin{tabular}{p{2.7cm} p{1.5cm} p{8.5cm}}
\toprule
\textbf{Library} & \textbf{Version} & \textbf{Role in the pipeline} \\
\midrule

\textbf{Raphtory} &
  0.17.0 &
  Stores the temporal multigraph and runs the Rust implementation of
  \citet{paranjape2017motifs}'s \texttt{local\_temporal\_three\_node\_motifs}
  algorithm in parallel. Its full 40-role output is checked before 32 roles are
  kept. \\

\textbf{DuckDB} &
  1.5.4 &
  Queries the prepared Parquet directly, computes the baseline, and sorts events
  for each ordering variant with bounded memory. \\

\textbf{PyArrow} &
  23.0.1 &
  Streams the full HuggingFace edge source, writes the prepared Parquet files,
  and carries the ordered batches from DuckDB to Raphtory. \\

\textbf{pandas} &
  3.0.2 &
  Holds the address-level targets, fingerprints, model matrices, predictions, and
  metrics. The primary HF and Real-CATS event tables are not materialised
  together in pandas; the much smaller BLTE table is. \\

\textbf{NumPy} &
  2.4.4 &
  Provides numerical arrays and the deterministic random numbers used by the
  models, the resampling, and the metrics. \\

\textbf{scikit-learn} &
  1.8.0 &
  Supplies the group-aware and stratified splitting, the random forest, the
  standardisation and logistic regression, and the metrics. \\

\textbf{XGBoost} &
  3.2.0 &
  Provides the gradient-boosted robustness model for the combined RQ1 view. \\

\textbf{matplotlib} &
  3.10.8 &
  Produces the figures written for the Results chapter. \\

\bottomrule
\end{tabular}
\end{table}

NetworkX and the HuggingFace \texttt{datasets} client are used only by the
exploratory application and by the data download, not by the locked pipeline.
SQLite, which comes with Python, provides the index used to de-duplicate the
Real-CATS histories.

\subsection{Project structure and entry points}

The repository keeps the exploratory and the evidence-producing code apart. The
interactive \texttt{main.py} is for development and inspection, while the thesis
results come only from \texttt{scripts/run\_results.py}, whose \texttt{prepare} and
\texttt{execute} subcommands call the separate \texttt{src/experiments} package.
The same entry point provides read-only \texttt{preflight} and \texttt{verify}
commands, execution \texttt{--resume}, and a coordinated
\texttt{execute-parallel} command for running the HF and Real-CATS roots on one
server.
Keeping the two apart stops exploratory shortcuts, such as row limits or changed
defaults, from leaking into the results. The main files are shown below, and
Figure~\ref{fig:pipeline_architecture} shows how they connect across the two
phases of a run.

\begin{lstlisting}[caption={Project layout (abbreviated).},
  label={lst:layout}]
bachelor_thesis/
main.py              CLI entry point, pipeline orchestrator
config.py            Exploratory defaults only
requirements.txt     Environment definition
requirements-results.txt Exact numerical Results environment
scripts/run_results.py  Locked Results prepare/execute CLI
src/
 experiments/
   constants.py     Locked constants and verified role catalogue
   prepare.py       Source conversion, cohorts, folds, protocol lock
   datasets.py      HuggingFace streaming and BLTE loading
   realcats.py      SQLite deduplication and layer export
   ordering.py      DuckDB external event ordering
   features.py      Baseline and fingerprint transformations
   motifs.py        Arrow-to-Raphtory loading and motif sweep
   modeling.py      Fixed estimators and development selection
   statistics.py    Metrics, permutations, bootstrap, BH correction
   splits.py        Group-aware and stratified partition construction
   rq_binary.py     RQ1/RQ2 execution and persistence
   rq_multiclass.py RQ3 execution and persistence
   secondary.py     Conditional BLTE and synthetic validation
   synthetic_fixtures.py  Deterministic motif-role and boundary fixtures
   execute.py       Verification and RQ orchestration
   execution_control.py Locks, attempts, threads, state transitions
   execution_artifacts.py Checkpoints, stages, result inventory
   coordination.py  Same-host HF/Real-CATS coordination
   verify.py        Read-only integrity and citability checks
   validation.py    Prepared-data and result-schema guards
   artifacts.py     Exclusive writes, hashes, runtime provenance
   resources.py     Read-only RAM/disk projection and admission
   telemetry.py     Fsynced JSONL timing and five-minute heartbeat
   types.py         Shared data contracts (dataclasses)
   visualization.py Chapter-ready PDF figures
 graph_construction.py  Exploratory NetworkX graph features
 motifs/            Exploratory alternative motif backends
data/raw/            Raw datasets (gitignored)
results/experiments/ Locked run directories
\end{lstlisting}

\begin{figure}[ht]
  \centering
  \ResGraphic{1.0}{01-pipeline}
  \caption{The locked Results pipeline. The preparation phase fixes the cohorts
  and writes the locked artifacts, and the execution phase then computes the
  motifs, trains the models, and runs the statistical tests. Italic labels name
  the source modules in \texttt{src/experiments}.}
  \label{fig:pipeline_architecture}
\end{figure}

\subsection{Engineering challenges}

The final design came out of an audit of the earlier exploratory pipeline against
what the three research questions actually need. The audit turned up both speed
bottlenecks and risks to the validity of the results, and the decisions below
address them.

\subsubsection{Locating the bottleneck}

The largest data problem is the size of the HuggingFace source, which is never
held in memory as a whole. The preparation phase scans the locally cached table in
small batches, measures how active each target is, and writes out only the
transfers that touch a selected address. This keeps the memory demand low, and it is
exact for the star and two-node motifs, though not for triangles, which is why
those are excluded (Section~\ref{sec:graph_construction}) rather than recorded as
zeros. Keeping the source locally removes download time and network failures, but
does not remove these complete local scans.

Timing the earlier pipeline showed that reading the data and fitting the models
were not the main cost. The dominant operation is Raphtory's local temporal motif
routine. The complete HF protocol requires seven windows under seven timestamp
orderings, and Real-CATS repeats the same grid within five layers. The graph is
therefore built once per ordering (and, for Real-CATS, per layer and ordering)
and reused across its seven windows, instead of being rebuilt for each window.

\subsubsection{Alternatives considered}

Three possible reductions were considered. First, counting only the motif roles
suggested by the scam examples would appear to save work. Raphtory~0.17.0,
however, exposes the 40-role counter as one operation. It cannot be instructed to
compute only selected roles. The pipeline consequently verifies the full output
and retains the 24 star and eight two-node roles. Avoiding the eight triangle
calculations would require a separate verified Rust implementation and was not
introduced during this study.

Second, replacing pandas with Spark would not accelerate the actual counter,
because the Raphtory graph still lives in one process. Pandas is used only after
aggregation, for address-level targets, fingerprints, model matrices,
predictions, and metrics. PyArrow performs bounded source scans and DuckDB
performs the baseline query and external event sorting. A distributed dataframe
rewrite would add another execution system without removing the single-process
motif bottleneck.

Third, reducing the transaction types was considered. The HF graph already
contains ordinary transaction edges only. For RQ3, in contrast, the separation
of normal, internal, ERC-20, ERC-721, and ERC-1155 histories is part of the
question: an exploit or rug pull may differ precisely in which layer carries its
activity. Dropping layers would change the fingerprint definition. The five
layers are therefore kept, but loaded sequentially so that only one layer graph
sets the Real-CATS peak memory demand.

\subsubsection{Long-run execution}

The workflow runs on a single machine, because the Raphtory graph has to live in
one process. A read-only preflight estimates graph expansion, DuckDB memory,
model matrices, checkpoint staging, sort spill, and a per-filesystem reserve.
This is an admission budget rather than an enforceable Raphtory memory limit. If
the budget does not fit, execution stops before creating an attempt instead of
quietly shrinking the locked cohort.

Each expensive motif call is committed independently, at ordering--window level
for HF and layer--ordering--window level for Real-CATS. A resumed execution
rechecks a canonical content hash and reuses valid commits; an invalid committed
checkpoint stops the run and is never replaced. HF and Real-CATS can also run in
parallel on one server through a coordinator that gives them disjoint thread
allocations and admits their combined RAM and disk demand. If one worker fails,
the other is allowed to finish. These changes do not reduce the number of motif
calculations in the scientific protocol, but they prevent an interruption late
in a week-long run from discarding earlier verified work.

\subsection{Reproducibility}

The study is built around a strict separation between fixing the experiment and
running it, so that no choice can be influenced by the test results. All settings live
in one place, \texttt{src/experiments/constants.py}, and are copied into the run's
protocol when preparation finishes. The main ones are shown below.

\begin{lstlisting}[language=Python,
  caption={Selected locked experiment constants.},
  label={lst:config}]
PROTOCOL_VERSION = 4
RANDOM_SEED = 42
CV_SEEDS = (42, 43, 44)
DELTAS = (300, 1800, 3600, 21600,
          86400, 259200, 604800)
N_PERMUTATIONS = N_BOOTSTRAPS = 10_000
MOTIF_ROLE_NAMES = STAR_ROLE_NAMES + TWO_NODE_ROLE_NAMES
HF_COHORT_ELIGIBILITY = {
    "minimum_incident_edges": 3,
    "activity_upper_quantile": 0.995,
}
\end{lstlisting}

When preparation finishes it writes \texttt{protocol.json}, which records the data
(the source revisions and a checksum of the source, the prepared files, the exact
target addresses, and the cohort and fold assignments), every model, window,
ordering, resampling, and seed value, and checksums of the code and of the
installed package and Python versions. A checksum is a short code computed from a
file that changes whenever the file changes, so it flags any later edit. From that
point the protocol and the prepared data are fixed. When a run is executed it
recomputes all of these checksums and refuses to produce a single feature or metric
if the code, the data, the versions, or the settings differ from what was locked. Fixed seeds make
the cohort sampling, the folds, and the resampling repeatable, and the
deterministic synthetic fixtures of Section~\ref{sec:data_acquisition} check that
the motif counter assigns the roles it should, including right at the edge of a
time window. The thesis profile fixes all settings above. A separate development
profile may reduce only the numbers of permutation and bootstrap repetitions and
is always marked non-citable.

Before execution, a read-only preflight checks the recorded RAM and disk budget
without creating a lock, log record, temporary directory, or timestamp change in
the run. Execution then creates an immutable
\texttt{execution\_definition.json}. This second definition separates operational
choices---the thread allocation, direct or coordinated mode, coordinator
identity, and concrete worker thread settings---from the analytical model
parameters in the protocol. A resumed run must use the same allocation. The
thread limits are set before importing NumPy, scikit-learn, or the native
libraries, which prevents numerical thread pools from silently exceeding that
allocation.

Exact equality is required for integer motif counts, address order, splits,
selected windows, and hypothesis decisions. Independently recomputed floating
probabilities and metrics are compared within a numerical tolerance because
BLAS, XGBoost, and Parquet encoders need not be byte-identical under different
thread allocations. Reusing an existing verified checkpoint or stage, in
contrast, is byte-preserving. PDF timestamps are removed so their hashes do not
change only because of generation time.

Scientific guards run before finalisation. They reject overlapping or duplicate
targets, campaign groups that cross a split, inconsistent folds, missing
classes, negative motif counts, non-finite features or metrics, duplicate
prediction rows, probabilities outside $[0,1]$, and multi-class probability
vectors that do not sum to one within numerical tolerance. Conflicting
duplicate labels in Real-CATS are rejected, while identical duplicates are
collapsed and counted in the preparation audit. All five layer files are
required even when one is empty. Checkpoint content is hashed from a fixed
Arrow representation with named columns, 64-bit integer counts, and addresses
sorted by their UTF-8 bytes, independently of its Parquet storage encoding.

\subsection{Result artifacts}

Each run has its own directory that is never overwritten, so results from
different runs cannot be mixed, and a run may only create a file that does not yet
exist. Prepared edges, fingerprints, and test-set predictions are stored as
Parquet, the protocol and the test summaries as JSON, the chapter tables as CSV,
and the figures as PDF. Output files are first fully written and synchronised in
a private attempt directory. Files are committed without replacement, and a
multi-file stage is published by an atomic directory rename after its manifest is
written last. The public \texttt{features}, \texttt{metrics}, and related
directories contain hard links to verified stage payloads, preserving the
filenames consumed by the Results chapter.

Preparation and execution have separate resume rules. During Real-CATS
preparation, the SQLite index records each completed address and can be resumed
only before the protocol is locked. During protocol-v4 execution, every attempt
has its own ID and scratch directory. HF commits one motif checkpoint for every
ordering and window; Real-CATS commits one for every layer, ordering, and
window. A new attempt may reuse earlier valid commits. A completed run rejects
both a fresh execution and a resume.

The integrity chain is acyclic. The locked protocol and prepared data identify
the execution definition; those identify every checkpoint and stage; the
result inventory hashes the committed stage manifests and public files; and the
execution manifest records the inventory hash. Each inventory and manifest has
a separate checksum file and does not include itself. The final
\texttt{verify} command checks this chain and the scientific schemas without
writing to the run. Ordinary verification distinguishes a valid non-citable
development or smoke run from corruption. Final thesis acceptance additionally
requires \texttt{--require-citable}.

\texttt{run\_log.jsonl} is the only appended file. Its writes are thread-safe
and synchronised to disk. A native operation that runs for several minutes
emits a heartbeat every five minutes with its elapsed time, allocated threads,
free disk, and latest committed checkpoint. The output layout is shown below.

\begin{lstlisting}[caption={Locked Results artifact layout.},
  label={lst:filenames}]
results/experiments/<run-id>/
  preparation_state.json
  protocol.json / protocol.sha256
  execution_definition.json / execution_definition.sha256
  attempts/<attempt-id>/
    started.json, completed.json or failed.json
    scratch/duckdb/
    staging/
  checkpoints/<checkpoint-id>/
    counts.parquet
    metadata.json
  stages/<stage-id>/
    payload/
    stage_manifest.json
  run_log.jsonl
  data/
  features/
  predictions/
  metrics/
  figures/
  tables/
  result_artifacts.json / result_artifacts.sha256
  execution_manifest.json / execution_manifest.sha256
\end{lstlisting}


Develop your hypotheses, describe your mental model for the research question.


\chapter{Results}
\label{ch:results}

This chapter reports the findings for the three research questions. RQ1 and RQ2
are evaluated on the two HuggingFace cohorts and RQ3 on the Real-CATS criminal
categories, all following the protocol fixed in
Section~\ref{sec:evaluation_protocol}. Every table and figure is produced directly
from the locked run artifacts.



\chapter{Discussion} What does it mean? Discuss the theoretical and practical implications of your findings as well as the limitations of your study.

\section{Theoretical Implications}\dots

\section{Limitations}\dots


\chapter{Conclusion \& Outlook} Repeat the problem and its relevance, as well as the contribution (plus quantitative results). Provide an outlook for further research steps.



Coordination patterns that flow primarily through smart contract layers (DeFi exploits,
 contract-mediated Ponzi schemes) are not captured by normal ETH transfers alone. Extending
 the pipeline to multi-layer temporal motif analysis across transaction types is left for future
 work.

 TODO: not entirely correct


%% ----------------
%% |   Appendix   |
%% ----------------
% IM Style: No additional blank page
%\cleardoublepage

%Use the appendix.tex file in the text folder to edit

\input{text/appendix} 




%% --------------------
%% |   Bibliography   |
%% --------------------
%%\cleardoublepage
\phantomsection
\addcontentsline{toc}{chapter}{\bibname}

% IM Style
\bibliographystyle{apalike}

%What research is referred to?

% Informatik-Style
%\iflanguage{english}
%{\bibliographystyle{IEEEtranSA}}	% english style
%{\bibliographystyle{babalpha-fl}}	% german style
												  
% Use IEEEtran for numeric references
%\bibliographystyle{IEEEtranSA})

\bibliography{thesis}


\end{document}

